\chapter{Проблемно-классификационный анализ задачи размещения объектов обслуживания населения с учетом изменчивости транспортной доступности}
\label{ch:problem-analysis}

\section{Размещение объектов обслуживания населения как задача территориального планирования}
\label{sec:service-placement-planning}

\section{Доступность как связанная характеристика застройки и транспорта}
\label{sec:accessibility-built-transport}

\section{Существующие методы оптимального размещения объектов обслуживания}
\label{sec:existing-location-methods}

\subsection{Модели set covering location problem}
\label{subsec:set-covering-location}

\subsection{Методы добавления новых объектов обслуживания}
\label{subsec:new-service-methods}

\subsection{Ограничения методов, рассматривающих спрос и транспортную сеть как константы}
\label{subsec:constant-demand-network}

\section{Изменчивость транспортной доступности как предмет исследования}
\label{sec:transport-accessibility-variability}

\subsection{Разрушение транспортных связей под воздействием внешних факторов}
\label{subsec:network-disruption}

\subsection{Улучшение транспортных связей как управляемая интервенция}
\label{subsec:network-improvement}

\subsection{Методы и алгоритмы поддержки принятия решений по обеспечению транспортной доступности}
\label{subsec:decision-support-accessibility}

\section{Выводы по главе}
\label{sec:chapter1-conclusions}
